\section{The Evaluation Crisis: Why the Numbers Cannot Be Trusted}
\label{sec:eval}

If the ladder tells us \emph{what} could be translated, the evaluation
problem tells us why we cannot yet know \emph{whether} it has been. This
section argues that the field's most quoted quantities---vendor accuracies
above 80\%, press-release superlatives---are, as presently defined,
\emph{unfalsifiable}, and that fixing this is a prerequisite for calling any
system a translator.

\subsection{The ground-truth problem}
Translation, in the human sense, is scored against a reference: a bilingual
corpus supplies the intended meaning, and metrics such as BLEU or human
adequacy judgments measure fidelity to it. Canine communication has no such
reference. There is no canonical ``intended message'' behind a bark, no
bilingual informant who can confirm that ``I want to play'' is what the dog
meant. The dog cannot report its intent, and---per
Section~\ref{sec:biology}---there may be no sentence-shaped intent to report.
Consequently, a bark-to-sentence system cannot be scored for translational
fidelity at all. What \emph{can} be scored is the upstream classifier: does
the model's emotion label agree with a human consensus label on a held-out
recording? That is a Rung-2 metric, and it is the only honest number in the
pipeline. Reporting it as a ``translation accuracy'' is a category mistake.

\subsection{What the vendor numbers actually are}
When a product advertises ``81.5\% accuracy,'' the defensible reading is:
on some internal test set, the predicted emotion class matched a
human-assigned class 81.5\% of the time. Three problems follow. First, the
test set, label protocol, class balance, and inter-annotator agreement are
rarely disclosed, so even the Rung-2 claim cannot be independently checked.
Second, the label is the \emph{human-perceived} emotion or the
\emph{eliciting context}, which is a proxy for the dog's state, not a
readout of it; the ELTE tradition is explicit about
this~\cite{pongracz2006acoustic,icassp2018dogbark}. Third, the rendered
sentence is produced \emph{after} classification by a language model
optimized for human resonance, so its fidelity is never measured at all.
The number thus travels down the pipeline as marketing while its meaning
stays fixed at the top.

\subsection{Unfalsifiability and the absent complaint}
A scientific claim must be falsifiable: there must be some observation that
could show it false. ``This bark means `I love you'\,'' admits no such
observation---the dog will not contradict it, and any subsequent behavior
can be retrofitted to the narrative. This is the structural reason the
consumer market cannot self-correct: there is no error signal. Compare a
weather forecast, which is scored against tomorrow's weather, or a
speech-to-text system, scored against a transcript. The pet translator has
no tomorrow and no transcript. Until a validation loop supplies an error
signal, its outputs are unfalsifiable interpretations, not measurements.

\subsection{Clever Hans, anthropomorphism, and the observer effect}
Two classic threats compound the problem. The \emph{Clever Hans} effect---%
the horse that appeared to do arithmetic but was reading its trainer's
unconscious cues---is the permanent hazard of any animal-communication
experiment: the human observer leaks the answer. Soundboard studies guard
against this with blind protocols and unfamiliar
button-pressers~\cite{bastos2024plos}; consumer apps, used by the devoted
owner in an uncontrolled home, do the opposite. \emph{Anthropomorphism} is
the second threat: attributing human propositional attitudes (``she says
she loves me'') to graded affective signals. The critical-review literature
on canine emotion perception warns that human observers project their own
categories onto the dog~\cite{sci2022explainable}, and a language model
trained on human text will amplify exactly this projection. The product
thus risks teaching owners to \emph{ignore} the dog's real body signals in
favor of a flattering sentence---a welfare harm we return to in
Section~\ref{sec:ethics}.

\subsection{What a valid evaluation would require}
A defensible evaluation program has four ingredients.
\begin{enumerate}[label=(\arabic*)]
  \item \textbf{Defined labels and disclosed protocols.} Report Rung-2
  metrics on a public, held-out corpus with inter-annotator agreement, class
  balance, and per-class confusion---the BEANS/BirdSet discipline applied to
  dogs~\cite{hagiwara2023beans}.
  \item \textbf{Calibration, not just accuracy.} Output probabilities that
  match observed frequencies, so a welfare tool can threshold responsibly.
  \item \textbf{Closed-loop validation.} Test inferred intent against the
  dog's subsequent behavior under blind conditions---the playback logic of
  Section~\ref{sec:pipeline}, adapted to the home.
  \item \textbf{Separate the sentence from the score.} Never report a
  generative rendering as evidence of fidelity; treat it as a user-interface
  layer, clearly labeled as interpretation.
\end{enumerate}
None of today's consumer products satisfies these. The research community
satisfies (1) and (2) in patches and (3) in the soundboard and cross-species
literature. The gap between the two is the evaluation crisis.

\subsection{A falsifiability criterion for the field}
We propose a simple demarcation: \emph{a canine communication system should
be called a translator only if its outputs are scored against an
independent, animal-grounded error signal.} By this criterion, no current
product qualifies; the soundboard paradigm qualifies at the word level;
and the rest is affect classification with a narrative interface. This is
not pedantry. It is the difference between a measurement and a story, and
the field's credibility---and the animals' welfare---depends on keeping them
apart.